\section{Introduction}
\label{sec:v86-introduction}
An action enters a clocked observation through a factor $T-x$. Unknown
selection changes its amplitude, normalization removes total mass,
and unknown readout channels conceal the identity of the component.
Recovering $x$ requires separating these three losses of information.
A fourth distinction is equally important: a detector that can be
rechosen independently at every site is not the same nuisance as one
shared over an observation domain. This paper studies the resulting
observation quotients, including their singular statistical boundary.

The first main result concerns spatial sharing. Suppose the latent odds
at two clocks are $r_j(z)=k_j(T_j-x(z))$, where the positive constants
$k_j$ are unknown and common to the sites, while the two full-rank
readout channels may vary with $z$. If the continuous action is not
constant on any open set, the observed matrix fields determine it and
both constants. Theorem~\ref{thm:v86-shared} proves this without
supplying labels or excluding isolated profile coincidences. The
proof identifies the exact obstruction to spatial calibration: a false
choice of clock constants allows directly matched action values at at
most one level and interchanged values at at most two levels. Four
nondegenerate calibration sites therefore determine the clock constants
by a finite calculation. On a simple surface, this gives a sharp
two-clock geometric threshold for a shared detector and constant unknown
reference. The one-clock lower bound uses positive full-rank channels
and distinct metrics, while keeping that reference and detector fixed.

The second main result describes the price of near rank deficiency.
At one site with binary channels, a constant relative detector and
three clocks, write $\rho=|\det U\det V|$. Theorem
\ref{thm:v86-product-modulus} proves, uniformly against all candidates
in the closed forward model, the action bound
\[
       |x-x'|\le C\min\{1,\|P-P'\|_*/\rho\}.
\]
The key step is to take the determinant after clearing the common pole
of the matrix curve. Its coefficient error has no extra inverse power
of the visible contrast. This yields honest confidence intervals with
locally adaptive expected length of order
$\min\{1,(\sqrt N\rho)^{-1}\}$ and a minimum-residual estimator
with squared risk of order $\min\{1,(N\rho^2)^{-1}\}$.
Theorems~\ref{thm:v86-adaptive} and \ref{thm:v86-singular-lower}
give matching uniform risk bounds and matching honest-length lower
bounds at explicit least-favourable parameters. Along the lower-bound
family both observed marginals remain exactly fixed; only a joint
interaction of size $\rho$ changes. Thus two equally weak channels of
contrast $\epsilon$ have transition $N\epsilon^4\asymp1$, rather
than an unspecified loss of conditioning. At zero contrast the family
is exactly indistinguishable, and honest intervals need not shrink.
The estimator does not receive a lower bound on $\rho$.

The third result is an intrinsic detector classification. On a compact
real Riemann surface, replace the logarithmic action curve by primitives
of normalized third-kind differentials. A finite-dimensional detector
space consists of primitives of meromorphic differentials with a fixed
pole divisor. Theorem~\ref{thm:v86-divisor} gives necessary and
sufficient tests involving integral residue charges and normalized
second-kind directions. The tests use periods as well as principal
parts. Theorem~\ref{thm:v86-divisor-budget} supplies the uniform
clock budget $\deg D+2g+4$. On the sphere this reduces exactly to
the rational polar budget; on positive-genus surfaces it gives genuinely
nonrational detector examples. Theorem~\ref{thm:v86-branched} further
retains the original factor $T-x$ under a finite meromorphic clock map,
allowing, for example, elliptic-integral detector primitives. The budget
is sufficient, not asserted
minimal for every meromorphic subspace. The sharper polynomial count
uses the additional structure of oriented intervals.

\subsection{The role of the observation quotient}
There are three different quotients in these conclusions. At a single
site the logarithmic action curve is taken modulo a detector function
space; its secants and tangents give exact fibres and local inversion.
For shared apparatus the nuisance consists of functions of the clock
index alone on the product of clock and spatial domains; spatial
contrasts and the residual permutation ambiguity determine a different
threshold. At finite sample size the relevant object is the modulus
of the full matrix observation map, after allowing its unknown channels.
The product theorem computes that modulus rather than inserting an
unknown condition number into a statistical bound.

These distinctions also delimit the geometric conclusions. Once an
experiment determines boundary distances, a separate geometric theorem
is needed to determine a metric. The surface rigidity input is the
Pestov--Uhlmann theorem \cite{V84PU}. The finite noisy metric statements
later in the paper use the specified local stability class and yield
confidence and stability certificates modulo boundary-fixing
diffeomorphisms. They do not claim an efficient inversion algorithm for
an arbitrary infinite-dimensional metric class, a new global
boundary-rigidity theorem, or a new partial-data geometric theorem.
The new sharp threshold concerns which selective observations determine
the geometric theorem's input.

\subsection{Relation to earlier work and organization}
Observable matrix pencils belong to the established identification
theory of latent models \cite{V83AMR,V83BJR}. We use that factorization
and resolve its remaining permutations by spatial sharing; the
factorization itself is not a novelty claim. Modulus-based recovery and
honest adaptation have classical precedents in
\cite{V86Donoho,V86CaiLow,V86Dufour}. Our statistical contribution is
the explicit uniform determinant-product modulus and its attaining
unknown-channel family, not a general two-point lower-bound method.
Normalized differentials and divisor degrees are classical tools
\cite{V86Forster,V86GKN}; our additional statement classifies which
unknown meromorphic detector subspaces preserve an action quotient.
Closest-predecessor comparisons are given again at the corresponding
theorems.

Sections~\ref{sec:v86-shared}--\ref{sec:v86-meromorphic} prove these
three results. The subsequent part retains the pointwise polynomial
and nonpolynomial inverses, rational examples, complete fibres,
finite-record covering reconstruction, geometric certificates, raw
acquisition bounds and exact lift constructions. Their observation
hypotheses are not changed by the shared-apparatus result. The expanded
edition additionally preserves the companion observation regimes and
historical derivations, with their own hypotheses. No conclusion from
one regime is silently used as an observation premise in another.
